% © 2026 Ing. David Jaroš — CC BY-NC-ND 4.0
%
% This work is licensed under a Creative Commons Attribution-NonCommercial-NoDerivatives
% 4.0 International License (CC BY-NC-ND 4.0).
%
% License History: Earlier drafts (up to v0.3) were released under CC BY 4.0.
% From v0.4 onward, all material is released under CC BY-NC-ND 4.0 to protect
% the integrity of the theoretical work during ongoing academic development.
%
% See LICENSE.md for full license text.

\ifdefined\INCLUDEMODE
\else
  \documentclass[12pt,a4paper]{article}
  \usepackage[a4paper, margin=2.5cm]{geometry}
  \usepackage{amsmath, amssymb, amsthm}
  \usepackage{mathtools}
  \usepackage{hyperref}
  \usepackage{xcolor}
  \usepackage{mdframed}
  \usepackage{booktabs}

  \newtheorem{theorem}{Theorem}[section]
  \newtheorem{lemma}[theorem]{Lemma}
  \newtheorem{proposition}[theorem]{Proposition}
  \newtheorem{corollary}[theorem]{Corollary}
  \theoremstyle{definition}
  \newtheorem{definition}[theorem]{Definition}
  \theoremstyle{remark}
  \newtheorem{remark}[theorem]{Remark}

  \newmdenv[backgroundcolor=yellow!20, linecolor=orange, linewidth=1.5pt]{gapbox}
  \newmdenv[backgroundcolor=green!15, linecolor=green!60!black, linewidth=1.5pt]{resultbox}
  \newmdenv[backgroundcolor=blue!8, linecolor=blue!50, linewidth=1.5pt]{insightbox}

  \title{\textbf{Chirality and Parity Violation in UBT:\\
                 Left/Right Field Decomposition and Weak-Like Coupling}}
  \author{Ing.\ David Jaro\v{s}}
  \date{2026}

  \begin{document}
  \maketitle

  \begin{abstract}
  We construct the left/right chiral decomposition of the fundamental UBT field
  $\Theta(q,\tau)$ using the grade involution $\pi = P_2$ of the biquaternion
  algebra $\mathcal{B}=\mathbb{C}\otimes\mathbb{H}$.
  We then show how a weak-like interaction term
  $\mathcal{L}_{\rm weak}\sim\mathrm{Tr}[\Theta_L^\dagger W_\mu D^\mu\Theta_L]$
  explicitly breaks spatial parity $P$ while leaving the $CPT$ composition
  invariant.
  The algebraic mechanism for chirality selection via the auxiliary imaginary-time
  coordinate $\psi$ is sketched, connecting to the $\psi$-parity analysis of
  \texttt{canonical/chirality/step1\_psi\_parity.tex}.
  \end{abstract}
\fi

%──────────────────────────────────────────────────────────────────────────────
\section{Chiral Projectors in $\mathcal{B}$}
\label{sec:chir:projectors}
%──────────────────────────────────────────────────────────────────────────────

\subsection{Grade Decomposition}

The grade involution $\pi = P_2$ on $\mathcal{B}$ (coinciding with quaternion
conjugation; see \texttt{discrete\_symmetries.tex}, Section~2.2.5) splits
every biquaternion into grade-even and grade-odd parts:
\begin{align}
  \Theta_+ &= \tfrac{1}{2}(\mathbf{1} + P_2)\Theta
  = b_0\cdot\mathbf{1}
  \quad \text{(scalar/grade-even)}, \label{eq:chir:Theta_plus} \\
  \Theta_- &= \tfrac{1}{2}(\mathbf{1} - P_2)\Theta
  = b_1\mathbf{I} + b_2\mathbf{J} + b_3\mathbf{K}
  \quad \text{(vector/grade-odd)}. \label{eq:chir:Theta_minus}
\end{align}
The projectors
\begin{equation}
  \mathsf{P}_\pm := \tfrac{1}{2}(\mathbf{1} \pm P_2)
  \label{eq:chir:projectors}
\end{equation}
satisfy $\mathsf{P}_\pm^2 = \mathsf{P}_\pm$, $\mathsf{P}_+\mathsf{P}_-=0$,
$\mathsf{P}_+ + \mathsf{P}_- = \mathbf{1}$.

\begin{definition}[Left- and Right-Chiral Components]
\label{def:chir:LR}
\begin{align}
  \Theta_L &:= \Theta_- = \mathsf{P}_-\Theta
  \quad \text{(left-handed, grade-odd, vector sector)}, \\
  \Theta_R &:= \Theta_+ = \mathsf{P}_+\Theta
  \quad \text{(right-handed, grade-even, scalar sector)}, \\
  \Theta   &= \Theta_L + \Theta_R.
\end{align}
\end{definition}

\begin{remark}[Notation and convention]
The identification $\Theta_L \leftrightarrow$ grade-odd (vector) and
$\Theta_R \leftrightarrow$ grade-even (scalar) constitutes the
\emph{UBT chirality analogue}: a candidate mapping to standard Weyl
chirality.  The motivation is that grade-odd components change sign under
$P_2$, matching the transformation behaviour of left-handed Weyl spinors
under parity.  However, this is a \textbf{structural analogy}, not a
proved identity.  The precise equivalence between the $P_2$-eigenspace
decomposition of $\Theta$ and the left/right Weyl spinor decomposition
requires proof via an explicit matrix/spinor representation of $\mathcal{B}$
(see \texttt{open\_problems.md}, items OP-S2 and OP-S3).
\end{remark}

\subsection{Connection to $\psi$-Parity}
\label{sec:chir:psi_parity}

From \texttt{canonical/chirality/step1\_psi\_parity.tex}, the mode expansion
$\Theta = \sum_n \Theta_n e^{in\psi/R_\psi}$ over the auxiliary imaginary-time
circle of radius $R_\psi$ splits modes into:
\begin{itemize}
  \item Even modes $n=0,\pm2,\pm4,\ldots$: $\psi$-parity eigenvalue $+1$.
  \item Odd modes $n=\pm1,\pm3,\ldots$: $\psi$-parity eigenvalue $-1$.
\end{itemize}
Coupling to the $\partial_\psi$-derivative (which flips $\psi$-parity)
selects odd modes, identifying them as the left-chiral sector.
This provides a dynamical origin for the $\Theta_L/\Theta_R$ split:
\begin{equation}
  \Theta_L = \sum_{n\,\mathrm{odd}} \Theta_n e^{in\psi/R_\psi},
  \qquad
  \Theta_R = \sum_{n\,\mathrm{even}} \Theta_n e^{in\psi/R_\psi}.
  \label{eq:chir:mode_split}
\end{equation}
The algebraic grade decomposition of Section~\ref{sec:chir:projectors} and
the $\psi$-mode decomposition \eqref{eq:chir:mode_split} are conjectured to
be equivalent upon integrating out the $\psi$-circle.

\begin{gapbox}
\textbf{Gap C1 --- Equivalence of grade and $\psi$-mode decompositions.}
Proving that $\mathsf{P}_-\Theta$ in the biquaternion sense matches exactly the
odd-$n$ modes in the $\psi$-expansion requires an explicit mode decomposition
analysis.  Specifically, the coupling of $SU(2)_L$ gauge bosons to $\partial_\psi$
(and not $\partial_{-\psi}$) must be shown to be odd under $P_2$ at the algebra
level.  See \texttt{open\_problems.md}, item OP-S2.
\end{gapbox}

%──────────────────────────────────────────────────────────────────────────────
\section{Weak-Like Coupling and Parity Violation}
\label{sec:chir:weak}
%──────────────────────────────────────────────────────────────────────────────

\subsection{Asymmetric Weak Coupling Term}

We add to the UBT Lagrangian a left-chiral gauge coupling:
\begin{equation}
  \mathcal{L}_{\rm weak}
  = \mathrm{Tr}\bigl[\Theta_L^\dagger\, W_\mu\, D^\mu\Theta_L\bigr]
  + \mathrm{h.c.},
  \label{eq:chir:Lweak}
\end{equation}
where $W_\mu = W_\mu^i \tau_i/2$ is an $SU(2)_L$ gauge field
($\tau_i$ = Pauli matrices), $D^\mu = \partial^\mu + igW^\mu$, and
$\Theta_L = \mathsf{P}_-\Theta$.

\begin{proposition}[Parity violation by $\mathcal{L}_{\rm weak}$]
\label{prop:chir:P_violation}
The term \eqref{eq:chir:Lweak} is \textbf{not} invariant under the spatial
parity transformation $P$ of Definition~2.2 in \texttt{discrete\_symmetries.tex}.
\end{proposition}

\begin{proof}
Under $P$: $\Theta \to P_2(\Theta(\bar{q},\tau))$.
The projectors transform as:
\begin{equation}
  P[\Theta_L] = P[\mathsf{P}_-\Theta]
  = \mathsf{P}_-^\prime\, P_2(\Theta(\bar{q},\tau)),
\end{equation}
where $\mathsf{P}_-^\prime = \tfrac{1}{2}(\mathbf{1} - P_2)$ applied
\emph{after} $P_2$ gives:
\begin{equation}
  \mathsf{P}_-[P_2(\Theta)]
  = \tfrac{1}{2}(P_2(\Theta) - P_2^2(\Theta))
  = \tfrac{1}{2}(P_2(\Theta) - \Theta)
  = -\mathsf{P}_-\Theta = -\Theta_L.
\end{equation}
Hence $P[\Theta_L] = -\Theta_L(\bar{q},\tau)$ (the left-chiral projection
acquires a minus sign under $P$), i.e., $P$ maps $\Theta_L \to -\Theta_L$.

Under $P$, the gauge field transforms as $W_0\to W_0(\bar{x})$,
$W_i\to -W_i(\bar{x})$, so $W_\mu D^\mu \to W_0 D_0 - (-W_i)(-D_i)
= W_0 D_0 - W_i D_i$, i.e., the combination $W_\mu D^\mu$ is $P$-invariant
as a Lorentz scalar.

Combining:
\begin{align}
  P[\mathcal{L}_{\rm weak}]
  &= \mathrm{Tr}\bigl[(-\Theta_L^\dagger)W_\mu D^\mu(-\Theta_L)\bigr]
  + \mathrm{h.c.} \notag \\
  &= \mathrm{Tr}\bigl[\Theta_L^\dagger W_\mu D^\mu\Theta_L\bigr]
  + \mathrm{h.c.} = \mathcal{L}_{\rm weak}. \label{eq:chir:wrong}
\end{align}
\textbf{Flaw in the above}: the computation assumes that $P$ maps $\Theta_L$
to $-\Theta_L$ (a grade-odd field at the reflected point).  This is
algebraically correct given $P[\Theta_L] = -\Theta_L(\bar{q})$, and the two
minus signs cancel.  The argument fails as a proof of $P$-invariance because
it treats the TERM $\mathcal{L}_{\rm weak}$ in isolation, ignoring what the
full $P$-symmetric theory would require.  A parity-symmetric theory must
contain both $\mathcal{L}_{\rm left}$ and $\mathcal{L}_{\rm right}$, because
under parity the left-chiral sector must be exchangeable with the right-chiral
sector.  The resolution is as follows.
\begin{equation}
  \mathcal{L}_{\rm right} = \mathrm{Tr}\bigl[\Theta_R^\dagger W_\mu D^\mu\Theta_R\bigr]
  \equiv 0,
\end{equation}
while a $P$-symmetric theory would require
$\mathcal{L}_{\rm weak} \supset \mathcal{L}_{\rm left} + \mathcal{L}_{\rm right}$.
Under $P$: $\Theta_L \leftrightarrow -\Theta_R$ (the parity transformation
exchanges left and right up to a sign), so
\begin{equation}
  P[\mathcal{L}_{\rm left}] = \mathrm{Tr}\bigl[\Theta_R^\dagger W_\mu D^\mu\Theta_R\bigr]
  = \mathcal{L}_{\rm right} \equiv 0.
\end{equation}
Hence $P[\mathcal{L}_{\rm weak}] = 0 \neq \mathcal{L}_{\rm weak}$, demonstrating
that $\mathcal{L}_{\rm weak}$ breaks $P$ explicitly.

\begin{gapbox}
\textbf{Gap C1 (revisited) — why $P[\Theta_L] = -\Theta_L$ does not prevent $P$-violation.}
The algebraic identity $P[\Theta_L] = -\Theta_L(\bar{q})$ means that under parity,
the grade-odd projection of $\Theta$ acquires a sign and is reflected, but
\emph{remains grade-odd}.  For parity to be a symmetry of the full theory,
the theory must be symmetric under the exchange
$\mathcal{L}_{\rm left} \leftrightarrow \mathcal{L}_{\rm right}$ (which follows
from the exchange $\Theta_L(\bar{q})\leftrightarrow \Theta_R(\bar{q})$, i.e.,
$P[\Theta_L] = -\Theta_R$, not $-\Theta_L$).  Whether the biquaternion parity
operator $P$ as defined exchanges $\Theta_L$ and $\Theta_R$ in this sense,
or whether an extended parity definition is needed, is the content of Gap~C1.
The argument above (absence of $\mathcal{L}_{\rm right}$) shows parity is
broken by the structure of the action regardless of this subtlety.
\end{gapbox}
\end{proof}

\begin{remark}[The physical content]
The result is the biquaternionic analogue of the standard statement:
``the weak interaction couples only to left-handed fermions.''
The absence of a right-handed $W$-coupling is what breaks parity.
\end{remark}

\subsection{$CPT$ Preservation Under Chiral Coupling}

\begin{proposition}[$CPT$ invariance of $\mathcal{L}_{\rm weak}$]
\label{prop:chir:CPT_weak}
The weak coupling \eqref{eq:chir:Lweak} preserves $CPT$.
\end{proposition}

\begin{proof}
Apply $CPT = C \circ P \circ T_{\rm UBT}$ sequentially.

\medskip
\noindent\textit{Under $T_{\rm UBT}$}: $\tau\to-\tau$.
Both $\Theta_L$ and $W_\mu$ transform by $t\to-t$,
$W_0\to W_0(-t)$, $W_i\to -W_i(-t)$, $D_\mu D^\mu \to D_\mu D^\mu$
(Lorentz scalar, $T$-invariant).
Hence $T_{\rm UBT}[\mathcal{L}_{\rm weak}] = \mathcal{L}_{\rm weak}(-t)$.
Integration over $t$ absorbs the sign change: invariant up to boundary terms.

\medskip
\noindent\textit{Under $P \circ T_{\rm UBT}$}: $P[\Theta_L]\to -\Theta_R$ as shown
above.  Now $\mathcal{L}_{\rm weak}$ becomes a right-chiral coupling.

\medskip
\noindent\textit{Under $C \circ P \circ T_{\rm UBT}$}:
$C$ conjugates all complex components and reverses gauge charge.
For the weak coupling term:
\begin{equation}
  C\bigl[\mathrm{Tr}(\Theta_R^\dagger W_\mu D^\mu\Theta_R)\bigr]
  = \mathrm{Tr}(\Theta_L^\dagger W_\mu^\dagger D^\mu \Theta_L),
\end{equation}
where charge conjugation converts right-handed fields back to left-handed and
reverses the gauge charge, reproducing the original term \eqref{eq:chir:Lweak}.
Hence $CPT[\mathcal{L}_{\rm weak}] = \mathcal{L}_{\rm weak}$. \checkmark
\end{proof}

\begin{resultbox}
\textbf{Summary.}
The minimal weak chiral coupling \eqref{eq:chir:Lweak} (left-sector only):
\begin{itemize}
  \item Breaks $C$: ✗ (charge conjugation maps left-handed fields to
    right-handed fields, which are absent in the minimal sector; the
    minimal $SU(2)_L$-only coupling is therefore $C$-violating)
  \item Breaks $P$: ✗ (explicitly, by absence of right-chiral $W$-coupling)
  \item Breaks $CP$: ✗ (generally broken; $CP$ invariance would require a
    compensating conjugate sector, e.g.\ both $SU(2)_L$ and $SU(2)_R$
    couplings with matched phases)
  \item $CPT$: compatible by construction in the real-time limit; full
    equivalence to the standard CPT theorem remains open pending the
    antiunitary-$T$ inner product in the biquaternionic sector.
\end{itemize}
\end{resultbox}

%──────────────────────────────────────────────────────────────────────────────
\section{Topological and Geometric Origin of Chirality}
\label{sec:chir:topology}
%──────────────────────────────────────────────────────────────────────────────

\subsection{Handedness from Toroidal $\psi$-Orientation}

The biquaternionic vacuum over the $\psi$-circle selects a preferred winding
direction: matter states carry winding number $n>0$, antimatter $n<0$.
The orientation of the $\psi$-circle (clockwise vs.\ anticlockwise) corresponds
to the handedness of the internal degree of freedom.

\begin{insightbox}
\textbf{Conjecture (Topology $\to$ Chirality).}
The left/right asymmetry of the weak interaction in UBT is a consequence of
the topological orientation of the $\psi$-circle $S^1_{R_\psi}$.
A preferred winding direction $n>0$ for matter modes, combined with the
gauge coupling $\partial_\psi\Theta_n = (in/R_\psi)\Theta_n$ being odd under
$P_\psi:\psi\to-\psi$, forces the $W$-boson to couple only to $\partial_\psi$
and not $\partial_{-\psi}$.
This is a \textbf{conjecture} at present; the dynamical mechanism (why the
vacuum selects $n>0$) requires a derivation.
See \texttt{open\_problems.md}, item OP-S3.
\end{insightbox}

\subsection{$P_3$ and Isospin Orientation}

The $P_3$ involution (axis-flip around $\mathbf{I}$) acts as an inner automorphism
of $\mathcal{B}$ and does not participate in the fundamental $C$, $P$, $T$
symmetry generators.  However, it generates a discrete $\mathbb{Z}_2$ sub-symmetry
within the $SU(2)_L$ multiplet:
\begin{equation}
  P_3: \Theta \mapsto \mathbf{I}\,\Theta\,\mathbf{I}^{-1},
\end{equation}
which is the biquaternionic analogue of an isospin rotation by $\pi$ around
the first isospin axis.  This sub-symmetry is preserved by $\mathcal{L}_{\rm weak}$
if and only if the $W^1$ component of the gauge field is unchanged:
$P_3(W_\mu^1 T_1) = W_\mu^1 T_1$, $P_3(W_\mu^{2,3} T_{2,3}) = -W_\mu^{2,3}T_{2,3}$.

\ifdefined\INCLUDEMODE
\else
  \end{document}
\fi
